% =============================================================================
\section{Motor Calibration}
\label{sec:motor}
% =============================================================================

\begin{center}
\noteto{Tiago}{Please review this whole motor calibration section -- check whether the description of the thrust-shaping, dead-band, ramp and FOC logic below is still accurate and consistent with the current code (\texttt{motor\_thrust2RPM}, VESC configs), and flag anything we should add, correct, or cut before this goes in the report.}
\end{center}

Thrust commands travel through three stages before becoming water flow, as sketched in Figure~\ref{fig:actuation}: (i)~\texttt{ControlLogic}/\texttt{VelocityController} compute a desired thrust in Newtons per side; (ii)~\texttt{motor\_thrust2RPM} converts that thrust command into a bounded, rate-limited ESC duty cycle; (iii)~the VESC driver forwards the duty cycle to the ESC, which runs closed-loop FOC commutation on the BLDC thruster motor.

\begin{figure}[H]
\centering
\begin{tikzpicture}[>=Stealth,node distance=6mm,every node/.style={font=\footnotesize}]
\tikzstyle{blk}=[draw,rounded corners,fill=orange!8,minimum width=3.0cm,minimum height=1.0cm,align=center]
\node[blk] (a) {Thrust command\\$T$ [N]};
\node[blk,right=16mm of a] (b) {motor\_thrust2RPM\\dead-band + ramp};
\node[blk,right=16mm of b] (c) {Duty cycle\\$[-0.9,\,0.9]$};
\node[blk,right=16mm of c] (d) {VESC (FOC)\\ERPM control};
\node[blk,right=16mm of d] (e) {Thruster\\thrust $T$ [N]};
\draw[->] (a)--(b); \draw[->] (b)--(c); \draw[->] (c)--(d); \draw[->] (d)--(e);
\end{tikzpicture}
\caption{Actuation chain from commanded thrust to delivered thrust.}
\label{fig:actuation}
\end{figure}

\subsection{Thrust-to-duty-cycle shaping (\texttt{motor\_thrust2RPM})}

The node first normalizes the requested thrust by the configured \texttt{max\_thrust}:
\begin{equation}
d_{\text{target}} = \frac{T_{\text{cmd}}}{T_{\max}}, \qquad d_{\text{target}} \in [-1,1].
\end{equation}
A dead-band suppresses commands too small to overcome static friction/cogging:
\begin{equation}
d_{\text{target}} = 0 \quad \text{if} \quad |d_{\text{target}}| < d_{\text{base}}, \qquad d_{\text{base}} = 0.08.
\end{equation}
The duty cycle is then rate-limited so the propeller never receives a step input. Three ramp modes are implemented; the deployed configuration uses \textbf{mode 2 (quadratic ease-out)}, in which the allowed rate of change grows quadratically with the remaining error, giving a soft start and a snappy finish:
\begin{align}
e_{\text{ramp}} &= \frac{|d_{\text{target}} - d_{\text{last}}|}{d_{\max}}, \quad \text{clamped to } [0,1] \\
\dot d_{\text{allowed}} &= \underbrace{0.2\, \dot d_{\max}}_{\text{base rate}} \;+\; \underbrace{20\, \dot d_{\max}}_{\text{explosive rate}} \cdot e_{\text{ramp}}^2, \\
\Delta d_{\max} &= \dot d_{\text{allowed}} \cdot \Delta t, \\
d_{\text{last}} &\leftarrow d_{\text{last}} + \operatorname{clamp}(d_{\text{target}}-d_{\text{last}},\,-\Delta d_{\max},\,\Delta d_{\max}),
\end{align}
where $\dot d_{\max}=$ \texttt{max\_duty\_cycle\_rate} $=0.9\,\text{s}^{-1}$. A zero-crossing (sign reversal of the command) triggers an instantaneous stop-before-reverse when \texttt{instant\_deceleration=true}, modeling the fact that water drag decelerates the hull quickly compared to thruster spin-up; same-direction deceleration is smoothed if \texttt{smooth\_same\_dir\_decel=true}. The result is clipped to $d\in[d_{\min},d_{\max}]=[-0.9,\,0.9]$ before being published as the duty-cycle command.

\begin{table}[H]
\centering
\caption{\texttt{motor\_thrust2RPM} parameters (identical for sim and real, \texttt{lily\_control/cfg/*.yaml}).}
\label{tab:t2rpm}
\begin{tabular}{@{}lcl@{}}
\toprule
Parameter & Value & Meaning \\
\midrule
\texttt{max\_duty\_cycle} / \texttt{min\_duty\_cycle} & $\pm 0.9$ & Saturation limits \\
\texttt{max\_duty\_cycle\_rate} & $0.9\,\text{s}^{-1}$ & Max slew of duty cycle \\
\texttt{ramp\_mode} & 2 (quadratic) & 0 none / 1 linear / 2 quadratic ease-out \\
\texttt{base\_duty\_cycle} & 0.08 & Static-friction dead-band \\
\texttt{startup\_thresh} & 0.10 & Reduced accel. below this target (mode 1) \\
\texttt{instant\_deceleration} & true & Snap to zero on direction reversal \\
\texttt{smooth\_same\_dir\_decel} & true & Ramp (not snap) when slowing without reversing \\
\texttt{loop\_rate} & 50\,Hz & Node update rate \\
\bottomrule
\end{tabular}
\end{table}

\subsection{VESC configuration and FOC commutation}

The duty-cycle command is sent over a serial link (\texttt{transport\_drivers/serial\_driver}) to a VESC-based ESC running in Field-Oriented Control mode. Table~\ref{tab:vesc} lists the deployed limits (\texttt{vesc/vesc\_driver/params/vesc\_config.yaml}).

\begin{table}[H]
\centering
\caption{VESC configuration limits.}
\label{tab:vesc}
\begin{tabular}{@{}lc@{}}
\toprule
Parameter & Value \\
\midrule
Serial port & \texttt{/dev/ttyACM0} \\
Duty cycle range & $[-1.0,\,1.0]$ (driver clamps to $\pm0.9$ upstream) \\
Motor current limit & $[0,\,100]$\,A \\
Speed limit (ERPM) & $\pm 23250$ \\
Brake current range & $[-20000,\,200000]$ \\
Servo output range & $[0.15,\,0.85]$ \\
\bottomrule
\end{tabular}
\end{table}

FOC commutation requires an in-water motor detection/calibration pass so that winding resistance, inductance and flux-linkage (and, if used, encoder/Hall offset) are identified under realistic load -- a submerged propeller presents a very different load curve than the motor spinning in air. \todored{Describe here the exact FOC detection procedure you ran in water (VESC ``Motor FOC $\rightarrow$ Detect'' wizard values: measured resistance, inductance, flux linkage $\lambda$, and any manual current/KV overrides), and the resulting current/torque constants ($K_t$, $K_v$) that were applied.}

\subsection{Thrust curve (ERPM $\rightarrow$ thrust)}
\label{sec:thrustcurve}

The relationship between commanded electrical RPM (ERPM) and delivered thrust is modeled with the standard quadratic propeller law (thrust $\propto n^2$ for fixed pitch, fixed advance ratio):
\begin{equation}
T(\text{erpm}) \;=\; T_{\max}\left(\frac{\text{erpm}}{\text{erpm}_{\max}}\right)^{2},
\label{eq:thrustcurve}
\end{equation}
with the reference values used to generate the calibration table (\texttt{vesc\_driver/scripts/motor\_curve.py}): $\text{erpm}_{\max}=23300$ (matching the VESC \texttt{speed\_max}$\,\approx 23250$ configured above) and $T_{\max}=200\,\text{N}$. This curve is what the FOC controller's closed-loop speed regulation and the thrust$\to$duty-cycle node above are calibrated against.

\begin{figure}[H]
\centering
\begin{tikzpicture}
\begin{scope}
\draw[->] (0,0) -- (7.2,0) node[right]{\footnotesize ERPM};
\draw[->] (0,0) -- (0,4.4) node[above]{\footnotesize Thrust [N]};
\draw[domain=0:7,samples=60,thick,blue] plot ({\x},{4*(\x/7)*(\x/7)});
\node at (3.5,-0.6) {\footnotesize $\text{erpm}_{\max}=23300$};
\node at (-1.1,4.1) {\footnotesize $T_{\max}=200$~N};
\draw[dashed,gray] (7,0) -- (7,4);
\draw[dashed,gray] (0,4) -- (7,4);
\end{scope}
\end{tikzpicture}
\caption{Nominal quadratic thrust curve, Eq.~\eqref{eq:thrustcurve}, used as the calibration reference.}
\label{fig:thrustcurve}
\end{figure}

\todored{This is the nominal/assumed quadratic model, not a fit to logged bench or in-water force-gauge data. If a real static-thrust-stand or in-water bollard-pull test was performed, insert the measured (erpm, thrust) data points here, refit Eq.~\eqref{eq:thrustcurve} (or a more general polynomial) against them, and report the fit residuals.}
